% GENERATED FILE. Edit canonical lessons, then run npm run generate:books.
% canonical-source-hash: fnv1a64:3fe36ccc06e21750
% canonical-lessons: PT-C16-ser, PT-C16-ser-roots, PT-C16-ser-vs-estar, PT-C16-ser-estar-meaning

\chapter{Ser and Estar}
\label{ch:pt-ser-and-estar}
\hlchaptermodality{16}

\begin{chapteropening}
\textbf{By the end of this chapter:} \emph{I can choose between ser and estar in Portuguese and change what an adjective means by switching them.}

Say Ela é bonita and Ela está bonita back to back, then A sopa é boa and a sopa está boa, naming the quality against the current condition each time.
\end{chapteropening}

\section[ser]{ser — “to be”}

\label{lesson:PT-C16-ser}



\begin{quote}
Portuguese has two verbs for “to be.” Begin with \textbf{ser}, used for identity and what something is; the choice against \emph{estar} comes later.
\end{quote}

\subsection*{You'll want to know: Present}
\noindent
\begin{tabularx}{\linewidth}{@{}>{\raggedright\arraybackslash}X>{\raggedright\arraybackslash}X@{}}
\toprule
\textbf{} & \textbf{} \\
\midrule
eu \textbf{sou} & nós \textbf{somos} \\
tu \textbf{és} & \textasciitilde{}\textasciitilde{}vós \textbf{sois}\textasciitilde{}\textasciitilde{} \\
ele/ela/você \textbf{é} & eles/elas/vocês \textbf{são} \\
\bottomrule
\end{tabularx}

\emph{Tu és} is live in European Portuguese. \textbf{Vós} survives mainly in northern dialect and church language; everyday “you all” is \textbf{vocês}, with \textbf{são}.

\subsection*{You'll want to know: Two past rows}
Preterite, the completed past from Chapter 15:

\begin{quote}
\textbf{fui, foste, foi, fomos, foram}
\end{quote}

Imperfect, for background or ongoing past:

\begin{quote}
\textbf{era, eras, era, éramos, eram}
\end{quote}

The three anchors \textbf{sou — fui — era} look unrelated. That is a clue that several historical pieces have been welded together, which the next lesson will unpack. For now, retrieve each row as a usable pattern.

\subsection*{Your turn}
Say these aloud:

\begin{itemize}
\raggedright
  \item “sou, és, é, somos, são”
  \item “Eu sou português. Ela é médica.”
  \item “fui, foste, foi, fomos, foram”
  \item “era, éramos, eram”
\end{itemize}

\begin{tcolorbox}[breakable,colback=teal!4,colframe=teal!35!black,title={Before you move on}]
Give the present of \emph{ser}. (\textbf{Sou, és, é, somos, são}.) Which form goes with \emph{vocês}? (\textbf{São}.) Give the preterite row. (\textbf{Fui, foste, foi, fomos, foram}.) Which imperfect form means “I was”? (\textbf{Era}.) Next: why those rows look so different.
\end{tcolorbox}
\section[sou / era / fui / ser]{ser — two verbs, three stems}

\label{lesson:PT-C16-ser-roots}



\begin{quote}
\textbf{Sou, era, fui, ser} do not resemble one another because the modern verb preserves several older roots.

\noindent
\begin{tabularx}{\linewidth}{@{}>{\raggedright\arraybackslash}X>{\raggedright\arraybackslash}X@{}}
\toprule
\textbf{Portuguese} & \textbf{source} \\
\midrule
\textbf{sou, és, é, somos, são} & Latin \emph{esse}, “to be” \\
\textbf{era, éramos, eram} & \emph{esse}’s imperfect \\
\textbf{fui, foi, foram} & \emph{esse}’s ancient perfect \emph{fuī} \\
infinitive \textbf{ser} & Latin \emph{sedēre}, “to sit” \\
\bottomrule
\end{tabularx}

That is \textbf{two Latin verbs} but \textbf{three stems}: \emph{esse} was already suppletive, using both its ordinary root and the older \emph{fuī} perfect. \emph{Sedēre} also gave English \textbf{sedentary} and \textbf{session}.
\end{quote}

\begin{cousinweb}
The entire preterite is shared with \textbf{ir}, “to go”:

\noindent
\begin{tabularx}{\linewidth}{@{}>{\raggedright\arraybackslash}X>{\raggedright\arraybackslash}X>{\raggedright\arraybackslash}X@{}}
\toprule
\textbf{person} & \textbf{ser} & \textbf{ir} \\
\midrule
eu & \textbf{fui} & \textbf{fui} \\
tu & \textbf{foste} & \textbf{foste} \\
ele & \textbf{foi} & \textbf{foi} \\
nós & \textbf{fomos} & \textbf{fomos} \\
eles & \textbf{foram} & \textbf{foram} \\
\bottomrule
\end{tabularx}

\textbf{Fui professor} means “I was a teacher”; \textbf{fui ao Brasil} means “I went to Brazil.” What follows disambiguates them. \emph{Fuī} already belonged to \emph{ser}’s ancestor; \textbf{ir} borrowed the row after Latin \emph{īre}’s own perfect wore away. Spanish made the same Iberian choice.
\end{cousinweb}

\subsection*{Your turn}
Say these aloud:

\begin{itemize}
\raggedright
  \item “ser $\leftarrow$ sedēre, to sit”
  \item “sou / era $\leftarrow$ esse”
  \item “fui professor; fui ao Brasil”
\end{itemize}

\emph{Twice through:} “Two verbs, three stems; \emph{fui} has two meanings.”

\begin{tcolorbox}[breakable,colback=teal!4,colframe=teal!35!black,title={Before you move on}]
Which Latin verb gives the infinitive? (\emph{\textbf{Sedēre}, “sit.”}) Which gives \emph{sou} and \emph{era}? (\emph{\textbf{Esse}.) What two verbs share }fui\emph{? (\textbf{Ser} and \textbf{ir}.) Which borrowed it? (\textbf{Ir}.) Next: put }ser\emph{ beside }estar\emph{.}
\end{tcolorbox}
\section[ser vs. estar]{ser vs. estar — the core choice}

\label{lesson:PT-C16-ser-vs-estar}



\begin{quote}
Chapter 2 used \textbf{Como está?} Now learn \emph{estar}’s forms and make the basic choice between Portuguese’s two verbs for “to be.”

\noindent
\begin{tabularx}{\linewidth}{@{}>{\raggedright\arraybackslash}X>{\raggedright\arraybackslash}X@{}}
\toprule
\textbf{} & \textbf{} \\
\midrule
eu \textbf{estou} & nós \textbf{estamos} \\
tu \textbf{estás} & \textasciitilde{}\textasciitilde{}vós \textbf{estais}\textasciitilde{}\textasciitilde{} \\
ele/ela/você \textbf{está} & eles/elas/vocês \textbf{estão} \\
\bottomrule
\end{tabularx}

Say \textbf{vocês estão} in ordinary speech. You may hear shortened \emph{tou} and \emph{tá}; recognise them, but write the full forms.
\end{quote}

\begin{grammarlens}[title={Grammar Lens: What something is; how it is}]
\begin{quote}
\textbf{ser} = what something is \textbf{estar} = its location or the condition it is in
\end{quote}

\noindent
\begin{tabularx}{\linewidth}{@{}>{\raggedright\arraybackslash}X>{\raggedright\arraybackslash}X@{}}
\toprule
\textbf{ser} & \textbf{estar} \\
\midrule
Ele \textbf{é} médico. — He is a doctor. & Ele \textbf{está} doente. — He is ill. \\
\textbf{Sou} português. — I’m Portuguese. & \textbf{Estou} em Lisboa. — I’m in Lisbon. \\
A neve \textbf{é} branca. — Snow is white. & O café \textbf{está} frio. — The coffee is cold. \\
\bottomrule
\end{tabularx}

Identity, origin, profession, and defining qualities use \textbf{ser}. Location, health, mood, weather, and resulting conditions use \textbf{estar}.

Do not reduce this to “permanent versus temporary”: \textbf{ele está morto}, “he is dead,” uses \emph{estar}. Death is not temporary; it is a condition someone has ended up in.
\end{grammarlens}

\begin{grammarlens}[title={What you've built: A mnemonic, not the cause}]
\emph{Ser} contains forms from Latin \emph{sedēre}, “sit”; \emph{estar} continues \emph{stāre}, “stand.” Picture \textbf{ser} as seated in identity and \textbf{estar} as standing in a condition. The history is a memory aid, not the cause: \emph{esse} supplied \emph{ser}’s identity meaning, while \emph{stāre} pulled toward place and resulting state.
\end{grammarlens}

\subsection*{Your turn}
Say these aloud:

\begin{itemize}
\raggedright
  \item “estou, estás, está, estamos, estão”
  \item “Sou português; estou em Lisboa.”
  \item “Ele é médico; ele está doente.”
\end{itemize}

\begin{tcolorbox}[breakable,colback=teal!4,colframe=teal!35!black,title={Before you move on}]
Which verb names identity? (\textbf{Ser}.) Which marks location or a resulting condition? (\textbf{Estar}.) How do you say “I’m in Lisbon”? (\textbf{Estou em Lisboa}.) Why does “he is dead” use \emph{estar}? (It is a \textbf{resulting condition}, so permanence is not the real test.) Next: use the choice to change meaning.
\end{tcolorbox}
\section[é bonita / está bonita]{é bonita / está bonita — a choice you can use}

\label{lesson:PT-C16-ser-estar-meaning}



\begin{quote}
\emph{Ser} and \emph{estar} are not competing answers where one is always a mistake. With many adjectives, swapping the verb deliberately changes meaning.

\noindent
\begin{tabularx}{\linewidth}{@{}>{\raggedright\arraybackslash}X>{\raggedright\arraybackslash}X@{}}
\toprule
\textbf{with ser} & \textbf{with estar} \\
\midrule
Ela \textbf{é} bonita. — She is beautiful. & Ela \textbf{está} bonita. — She looks beautiful now. \\
A sopa \textbf{é} boa. — It is a good soup. & A sopa \textbf{está} boa. — This bowl tastes good. \\
\bottomrule
\end{tabularx}

\textbf{Ser} presents the quality as part of what the subject is. \textbf{Estar} presents the quality as the current condition or impression. Context, not a mechanical “permanent” rule, decides what you mean.
\end{quote}

\begin{cousinweb}
Portuguese and Spanish maintain this full two-verb choice. Italian keeps \emph{essere} and \emph{stare} but lets them overlap differently; French has one \emph{être} that absorbed old \emph{stāre} pieces. German \emph{sein} is a separate Germanic cousin.

\noindent
\begin{tabularx}{\linewidth}{@{}>{\raggedright\arraybackslash}X>{\raggedright\arraybackslash}X@{}}
\toprule
\textbf{language} & \textbf{broad result} \\
\midrule
Portuguese / Spanish & two verbs, chosen by meaning \\
Italian & two verbs with overlap \\
French & one fused verb \\
German & one verb built from Germanic roots \\
\bottomrule
\end{tabularx}

The Romance rows started with the same Latin material and reorganised it in three ways. That is why the Portuguese choice is systematic but not universal.
\end{cousinweb}

\subsection*{Your turn}
Say these aloud:

\begin{itemize}
\raggedright
  \item “Ela é bonita” — a quality
  \item “Ela está bonita” — right now
  \item “A sopa é boa; a sopa está boa.”
\end{itemize}

\emph{Twice through:} “Ser: what it is. Estar: the condition it is in.”

\begin{tcolorbox}[breakable,colback=teal!4,colframe=teal!35!black,title={Before you move on}]
What changes between \emph{ela é bonita} and \emph{ela está bonita}? (A defining quality versus a \textbf{current impression}.) Is \emph{estar} automatically a mistake with \emph{bonita}? (\textbf{No}; it says something different.) Which Romance languages keep the full two-verb choice? (\textbf{Portuguese and Spanish}.)
\end{tcolorbox}
